\chapter{Conclusion} \label{ch:conclusion}

\section{Answering the Research Questions}

This chapter returns to the three research questions and summarizes what the
experiments show. The central finding is that constrained decoding converts
tool calling from a format-compliance problem, which is fixable structurally at
inference time, into a problem of semantic argument-value error. That residual
error can be reduced by restoring schema detail or, more modestly, by
format-aligned training. The answers below quantify these effects and then
translate them into deployment guidelines and future work.

\paragraph{RQ1.} \textit{What is the performance gap between unoptimized small
language models (0.5B--7B parameters) and frontier models on tool-calling benchmarks?}

Across all four Qwen~2.5 Instruct sizes evaluated on
BFCL~v4 \texttt{simple\_python}, the model-agnostic baseline achieves
1.5--4.75\% under the strict whole-completion parser, with no consistent
trend across model scale. The dominant failure mode is format non-compliance:
394 of 400 7B base outputs fail the strict parse rather than producing a
semantically wrong structured call, and the same pattern holds at all sizes.
This gap quantifies the cost of a naive generic integration and strict parse
contract rather than an intrinsic capability deficit: the same 7B run reaches
62.00\% when the first complete JSON object is accepted, and the B-template
control, which prompts the same model through its native tool-calling template
with free generation and full schemas, reaches 96.00\% (384/400) on the same
test cases.
Frontier models range from 78.75\% (GPT-5-mini) to 97.75\%
(Claude Sonnet~4.5) on the same category, with flagship models at
91.00\% or above~\cite{bfcl2024}; the B-template control falls inside
that flagship band.
Constrained decoding closes the strict-parser integration gap at every size
without modifying model weights, raising accuracy to 51.50\% (0.5B),
62.25\% (1.5B), 64.75\% (3B), and 72.75\% (7B), while guaranteeing that the
outputs are structurally parseable. At 7B, the same-parser gain over the
lenient unguided baseline is 10.75~pp.
At 7B, the residual gap to the best frontier model is approximately
25~pp, and 6.00~pp to the weakest frontier entry (GPT-5-mini,
78.75\%).
The residual gap is in argument value accuracy, where frontier models
and the native-template configuration share an advantage attributable
to prompts that carry the full schema detail and to stronger learned
associations for specific numeric, string, and format conventions.

\paragraph{RQ2.} \textit{What is the marginal contribution of each optimization
technique (constrained decoding, schema enrichment, prompt engineering,
chain-of-thought prompting, retrieval-augmented generation, quantization, and
LoRA fine-tuning) when applied cumulatively to
small models for tool calling?}

Constrained decoding is the dominant optimization technique because it gives a
model-agnostic structural guarantee: every output conforms to the required
schema, without training and without generating additional tokens. Relative to
the strict whole-completion parser floor, this corresponds to a $+48$--$+71$~pp
recovery across the 0.5B to 7B size range. At 7B, where the same unguided run
is also scored with the lenient first-object parser, the same-parser accuracy
gain is $+$10.75~pp (62.00\% to 72.75\%).

AWQ INT4 quantization reduces model memory by 63.5\% (14.25~GiB to
5.20~GiB at 7B) and startup time by 83.5\% (212~s to 35~s) at a
cost of at most 0.5~pp above 1.5B, within run-to-run
variance; at 0.5B the penalty is 3.75~pp.

Few-shot prompting (PE) produces a size-dependent effect: accuracy
improves by 2.0--2.5~pp at 0.5B--3B but decreases by 2.5~pp at
7B, indicating that the direction of
effect is not consistent across model scale.
Chain-of-thought reasoning (CD+Q+ITC) reduces accuracy at every
model size by approximately 6~pp at increased per-case inference time from the additional reasoning
tokens.
Retrieval-augmented generation (CD+Q+RAG) reduces accuracy at every
model size, with decreases of 16--26~pp relative to CD+Q, despite
a retrieval recall@5 of 97.2\%; the loss is attributable to function
disambiguation failure rather than retrieval quality.
The disambiguation failure does not improve with model scale: the
7B model is no more reliable than the 1.5B model at distinguishing
semantically similar candidate functions.

LoRA fine-tuning on
\texttt{xlam-function-calling-60k}~\cite{zhang2024xlam}
produced a 3~pp regression under CD+FT at 7B, attributed to a
format mismatch between the training distribution and the evaluation
pipeline.
The format mismatch hypothesis is confirmed by CD+FT-aligned, which
reformats training outputs to match the inference pipeline exactly.
CD+FT-aligned improves accuracy at every size: $+7.75$~pp at 0.5B,
$+3.75$~pp at 1.5B, $+2.00$~pp at 3B, and $+4.00$~pp at 7B
(76.75\%).
CD+Q+FT-aligned applies AWQ INT4 quantization to the format-aligned
LoRA-merged model and achieves 74.25\% (297/400) at 7B on the same
5.20~GiB footprint.
The quantization penalty relative to CD+FT-aligned is 2.5~pp at
7B, larger than the 0.5~pp penalty observed when quantizing the
base model, suggesting that AWQ interacts more destructively with
fine-tuned weight deltas when the calibration distribution does not
match the fine-tuning domain.
Paired McNemar tests place the 7B fine-tuning gain at the edge of
significance ($p = 0.044$, borderline against the observed
run-to-run spread) and the quantized-stack gain within noise
($p = 0.38$). These p-values are diagnostic rather than family-wise corrected
across the full ablation grid, so the consistent positive direction across all
four model sizes is the primary evidence for the format-alignment effect.

Config~CD+schema, which enriches the Stage~2 prompt with parameter
descriptions, enumeration values, and defaults while keeping the
FSM structural constraint unchanged, raises 7B accuracy from
72.75\% to 89.00\% ($+$16.25~pp, McNemar $p < 0.00001$;
Section~\ref{sec:cd-schema}).
This is the strongest no-training gain observed in this evaluation
and confirms that the dominant component of the residual after plain
CD is schema information loss rather than the model's learned value
priors.

The finding for RQ2 is that the residual failure rate after
constrained decoding is not addressed by generic prompt-level
additions at any model size; format-aligned training reduces it
modestly, and schema-enriched prompting reduces it substantially by
restoring the parameter context that the model-agnostic prompt
discards.

\paragraph{RQ3.} \textit{Can an optimized small language model configuration
provide an upper-bound local tier for cascade-style single-call tool calling on
consumer hardware?}

Under the scope defined in Chapter~\ref{ch:introduction}, the answer is yes as
an upper-bound local-tier result for single-call, schema-constrained tool
calling on consumer hardware. It is not a claim of end-to-end cascade
viability: the semantic router the cascade requires is not implemented, and
constrained decoding removes the structural validity signal that would otherwise
drive routing, so the reported traffic fractions are upper bounds
(Section~\ref{sec:practical-implications}). Config~CD+schema achieves
89.00\% at 7B (Section~\ref{sec:cd-schema}), which is 10.25~pp above
GPT-5-mini (78.75\%) and 2.00~pp below GPT-4.1 (91.00\%) on BFCL~v4 Python
Simple AST. It is the first constrained configuration in this evaluation to
enter the lower frontier range on that category while retaining structural
validity guarantees and requiring no model-weight updates. This comparison is
indicative because frontier scores are taken from the public leaderboard rather
than re-run under the local harness (Section~\ref{sec:frontier-comparison}).

The result is deliberately narrow. It does not establish broad production-agent
viability. On the multi-candidate and parallel BFCL categories, the same model
remains below the frontier band (Section~\ref{sec:bfcl-extended}), and on
$\tau$-bench retail the pass rate remains 4.35\% for CD, showing that
single-call gains do not transfer to multi-turn task completion. The B-template
control (96.00\%) also shows that, for a single frozen Qwen deployment, the
native tool-calling template is the stronger local baseline. The constrained
pipeline is justified when model-agnostic structural guarantees are required
across heterogeneous models and schemas.

Against the cascade criterion (a single consumer GPU and a local tier handling
$p \geq 0.70$ of single-call traffic), five configurations qualify.
CD+schema achieves 89.00\% at full bfloat16 precision (14.25~GiB), fitting on a
consumer NVIDIA RTX~4090 (24~GiB) alongside substantial key-value cache
capacity. CD+Q achieves 72.25\% on 5.20~GiB with no training required.
CD+Q+FT-aligned achieves 74.25\% on the same 5.20~GiB footprint, while
CD+FT-aligned achieves 76.75\% and CD 72.75\% at full bfloat16 precision
(14.25~GiB). The cascade interpretation presumes that residual semantic
failures are handled by escalation to a frontier model; the semantic router
needed to do this is not implemented in this thesis.

\begin{table}[ht]
\centering
\caption{Main thesis claims, supporting evidence, and scope.}
\label{tab:claim-evidence-scope}
\begin{tabular}{p{3.6cm} p{5.2cm} p{4.2cm}}
\toprule
Claim & Evidence & Scope \\
\midrule
Constrained decoding removes the dominant structural failure mode &
Config~B scores 1.50\%; Config~CD scores 72.75\%, with 100\% valid function selection on \texttt{simple\_python} &
BFCL~v4 \texttt{simple\_python}; model-agnostic Qwen~2.5 pipeline \\
Schema detail explains much of the residual after CD &
Config~CD+schema raises accuracy from 72.75\% to 89.00\% ($+$16.25~pp, McNemar $p < 0.00001$) &
7B model on BFCL~v4 Python Simple AST \\
Quantization is a practical deployment step above 1.5B &
AWQ INT4 reduces 7B memory from 14.25~GiB to 5.20~GiB with a 0.5~pp accuracy change &
Qwen~2.5 AWQ variants; single-call BFCL evaluation \\
Format-aligned fine-tuning helps, but less than schema enrichment &
CD+FT-aligned reaches 76.75\% at 7B and improves all four Qwen sizes, while CD+schema reaches 89.00\% without training &
Format-aligned xlam LoRA; BFCL single-call setting \\
Single-call gains do not solve multi-turn agency &
CD reaches 72.75\% on BFCL \texttt{simple\_python} but 4.35\% pass rate on $\tau$-bench retail &
$\tau$-bench retail with a 7B user simulator \\
\bottomrule
\end{tabular}
\end{table}

\section{Practical Deployment Guidelines}

The results suggest a decision procedure for teams deploying SLMs in
single-call tool-calling systems. The procedure is ordered by practical return:
use the strongest available integration first, restore schema detail before
training, and only then pay the cost of fine-tuning. The cascade cost model that
quantifies the saving from each step is developed in
Section~\ref{sec:practical-implications}.

\textbf{Step~1: Benchmark the native template and decide whether model-agnostic
guarantees are required.}
When the deployed model ships a tool-calling chat template, that template is the
zero-cost reference point and should be benchmarked first; in this evaluation it
reaches 96.00\% with free generation (B-template). For a single frozen Qwen
deployment, this is the stronger local baseline. Constrained decoding should be
used when structural validity must be guaranteed independently of the model's
native template, or when the deployment must support heterogeneous models and
schemas under one integration.

\textbf{Step~2: Apply constrained decoding with full schema detail.}
The highest-return no-training configuration is CD+schema. Carrying parameter
descriptions, enumeration values, and default values into the argument
extraction prompt raises 7B accuracy from 72.75\% to 89.00\% ($+$16.25~pp,
$p < 0.00001$) while preserving the structural guarantee of guided decoding.
This should be the default constrained setup whenever tool schemas contain rich
metadata. Plain CD remains useful when schemas are sparse or when only type
constraints are available.

\textbf{Step~3: Quantize when memory or cold-start latency matter.}
AWQ INT4 quantization reduces 7B model memory from 14.25~GiB to 5.20~GiB and
startup time from 212~s to 35~s, at a cost of 0.5~pp for the base constrained
model. This cost is within run-to-run variance at 7B, making quantization the
recommended deployment step when hardware capacity or load time is binding.

\textbf{Step~4: Treat generic prompting additions as low priority.}
Few-shot prompting and chain-of-thought reasoning both regress at 7B in this
evaluation: $-2.5$~pp for few-shot and $-6.75$~pp for chain-of-thought, with
additional token cost for the reasoning trace. These methods do not address the
residual failure mode after constrained decoding, which is semantic convention
matching rather than format learning.

\textbf{Step~5: Fine-tune only with deployment-aligned data.}
Format-aligned LoRA fine-tuning improves on plain CD, reaching 76.75\% at 7B,
but it is weaker than schema enrichment and depends on exact alignment between
the training output format and the inference protocol. Fine-tuning should be
reserved for cases where schema enrichment is insufficient and training examples
can be generated from the target tool schemas and argument conventions. Applying
the merged adapter to the quantized model yields 74.25\% on the same 5.20~GiB
footprint.

\textbf{Deployment boundary.}
These guidelines apply to single-call, schema-constrained scenarios analogous
to BFCL~v4 \texttt{simple\_python}. They do not establish reliable multi-turn
agentic execution. The $\tau$-bench result (4.35\% pass rate on retail) shows
that applications requiring multi-step tool use need additional planning,
state-tracking, and recovery mechanisms beyond the techniques evaluated here.

\section{Future Work} \label{sec:future-work}

The $\tau$-bench result (4.35\% pass rate on retail) establishes that
single-call accuracy does not predict multi-step agentic performance.
Improving multi-step reliability is the most important open direction.
The failure mode on $\tau$-bench is not argument formatting (constrained
decoding handles that) but multi-turn planning and error recovery across
5--15 tool calls. Techniques such as process reward models, chain-of-thought
with tool-observation grounding, and task-specific fine-tuning on multi-turn
trajectories are natural next steps~\cite{taubench2024}.
A controlled ablation isolating the contribution of each technique to
$\tau$-bench pass rate would determine which combination most effectively
closes the gap between the 4.35\% result observed here and the pass rates
reported for frontier models on the same benchmark.
The $\tau$-bench airline domain offers a second evaluation environment that
would indicate whether the retail result generalises across task structures
and tool catalog sizes.

Several extensions would complete the single-call evaluation grid. Config~CD+Q+FT-aligned is evaluated at 7B only; extending it to 0.5B, 1.5B, and 3B would show whether the 2.5~pp quantization penalty on the fine-tuned model increases at smaller scales. The no-guided isolation variants are also reported only at 7B; running them across the size range would test whether the unguided floor and the ordering of prompt-level techniques remain stable as capacity decreases.

A smaller controlled ablation would also sharpen the central positive result. Config~CD+schema injects three kinds of information at once: natural-language parameter descriptions, enumerated value constraints, and default values. The 16.25~pp gain over plain CD is therefore an aggregate over all three. Adding them one at a time, on the same 400-case split and with no new model runs beyond re-prompting, would attribute the gain to its components and indicate which class of schema metadata is worth the cost of authoring for a target tool catalogue.

RAG-based tool selection requires a stronger disambiguation signal than top-5 retrieval provides. The retrieval component itself is not the bottleneck: recall@5 is 97.2\% in Config~CD+Q+RAG. The failure occurs when the model must choose among semantically similar retrieved candidates. Future work should therefore target the selection step directly, for example with a smaller $k$, a learned re-ranker, or fine-tuning on tool-selection examples from the target catalogue.

The harder BFCL categories point to a second direction. The joint multi-call schema lifts \texttt{parallel} from a structural 0.0\% under the original single-extraction design to 79.0\% at 7B, confirming that the original limitation was the output schema. The remaining gap on \texttt{parallel\_multiple}, where format-aligned fine-tuning lowers accuracy relative to plain CD, suggests that future training data should target simultaneous multi-call generation rather than JSON format alignment alone.

Training data alignment remains the binding constraint on LoRA fine-tuning. A production fine-tuning run should generate examples from the target deployment's own tool schemas and argument conventions, rather than relying only on a general-purpose corpus such as xlam. A systematic comparison of synthetic, domain-matched, and human-curated examples would show how much of the residual semantic error can be closed by data quality alone.

Finally, the cross-family evidence in this thesis is a targeted external-validity check rather than a full replication. Config~CD and CD+FT-aligned are evaluated across contrast checkpoints, but quantized and agentic configurations remain concentrated on Qwen~2.5. Extending AWQ INT4, quantized fine-tuned variants, and $\tau$-bench evaluations across the same contrast families would test whether the deployment conclusions hold beyond the primary model family.

Beyond the specific extensions above, the ordering of results points to a broader question for how tool-calling agents are built. The single most effective intervention measured here is not a larger model, more training data, or an extra reasoning step, but the restoration of schema detail that was already present in the tool definitions and had simply been dropped from the prompt. A no-training, prompt-level change outperformed format-aligned fine-tuning at a fraction of the engineering cost. This suggests that effort spent making tool interfaces machine-legible (complete parameter descriptions, explicit enumerations, and stated defaults) may currently return more reliability per unit of work than further training of the model that consumes them. As tool ecosystems converge on shared interface standards, the schema becomes a shared leverage point that every model behind it inherits, whereas a fine-tuning gain is, as the xLAM result shows, specific to one model and one output convention. The evidence here is from a single benchmark category and should not be over-generalised, but it is at least suggestive that the field has invested heavily in adapting models to tools while under-investing in making tools describe themselves well to models.
